\chapter{Atividade 5 -- Contagem de números primos com OpenMP}
\label{chap:atvd05}

\section{Enunciado da tarefa}
\label{sec:atvd05-enunciado}

\begin{quote}
Implemente um programa em C que conte quantos números primos existem entre 2 e um valor máximo \(n\).

Depois, paralelize o laço principal utilizando a diretiva
\texttt{\#pragma omp parallel for}. A lógica original do algoritmo não deve ser alterada.

Compare o tempo de execução, os resultados da versão sequencial e os resultados da versão paralela utilizando diferentes valores de \(n\). Observe possíveis diferenças no resultado e no desempenho.

Discuta também os desafios iniciais da programação paralela, considerando especialmente correção e distribuição de carga. Para a medição do tempo, utilize uma abordagem baseada em \texttt{omp\_get\_wtime()}.
\end{quote}


% ============================================================
\section{Objetivos e requisitos}
\label{sec:atvd05-objetivos}

A atividade teve como objetivo investigar a transformação de um algoritmo
sequencial simples em uma implementação paralela baseada em OpenMP,
mantendo inalterada a lógica fundamental da computação.

O problema adotado consiste em determinar a quantidade de números primos
existentes no intervalo

\[
[2,n],
\]

para diferentes valores do limite superior \(n\).

Os objetivos específicos foram:

\begin{itemize}
    \item implementar uma versão sequencial da contagem de números primos;
    
    \item paralelizar o laço externo de candidatos utilizando
    \texttt{\#pragma omp parallel for};
    
    \item observar o efeito de uma atualização concorrente não sincronizada
    sobre um contador compartilhado;
    
    \item implementar uma versão paralela correta utilizando a cláusula
    \texttt{reduction};
    
    \item verificar a equivalência funcional entre a versão sequencial e a
    versão paralela correta;
    
    \item medir o tempo de execução utilizando \texttt{omp\_get\_wtime()};
    
    \item avaliar diferentes tamanhos do problema;
    
    \item calcular o speedup obtido pela paralelização;
    
    \item discutir overhead de paralelização, condição de corrida e
    distribuição de carga.
\end{itemize}

O número de threads foi mantido fixo em quatro durante todo o experimento.
O objetivo, portanto, não foi estudar escalabilidade em função da quantidade
de threads, mas isolar o efeito da paralelização do laço principal.


% ============================================================
\section{Fundamentação conceitual}
\label{sec:atvd05-fundamentacao}

\subsection{Teste de primalidade}

Um número inteiro \(x > 1\) é primo quando seus únicos divisores positivos
são 1 e o próprio \(x\).

Uma forma direta de verificar a primalidade consiste em testar possíveis
divisores. Não é necessário testar divisores superiores a \(\sqrt{x}\),
pois, caso \(x\) possua um fator maior que sua raiz quadrada, necessariamente
possuirá também um fator correspondente menor que ela.

Na implementação utilizada nesta atividade, após o tratamento do número 2,
todos os demais números pares são descartados imediatamente. Para números
ímpares, são testados apenas divisores ímpares.

A condição do laço foi expressa como

\[
d \leq \frac{x}{d},
\]

em vez de

\[
d^2 \leq x.
\]

As duas expressões são equivalentes no domínio considerado, mas a primeira
evita a necessidade de calcular explicitamente a raiz quadrada e evita uma
multiplicação que poderia causar overflow para valores suficientemente
grandes.


\subsection{Paralelismo de dados}

A classificação de um candidato como primo ou não primo é independente da
classificação dos demais candidatos.

Assim,

\[
\text{primo}(i)
\]

pode ser calculado independentemente de

\[
\text{primo}(j),
\qquad i \neq j.
\]

Essa característica torna o laço externo da contagem adequado ao
paralelismo de dados.

Conceitualmente, o espaço de candidatos pode ser distribuído entre as
threads:

\begin{verbatim}
Thread 0  ---> parte dos candidatos
Thread 1  ---> parte dos candidatos
Thread 2  ---> parte dos candidatos
Thread 3  ---> parte dos candidatos
\end{verbatim}

Cada thread executa a mesma função de primalidade sobre diferentes elementos
do intervalo.


\subsection{Diretiva \texttt{parallel for}}

OpenMP permite distribuir automaticamente as iterações de um laço por uma
equipe de threads por meio da diretiva

\begin{verbatim}
#pragma omp parallel for
\end{verbatim}

A diretiva combina a criação de uma região paralela com a distribuição do
espaço de iterações do laço.

Nesta atividade, apenas o laço que percorre os candidatos entre 2 e \(n\)
foi paralelizado.


\subsection{Laços aninhados e granularidade do paralelismo}

O algoritmo apresenta naturalmente dois níveis de repetição:

\begin{enumerate}
    \item um laço externo que percorre os candidatos;
    \item um laço interno que testa possíveis divisores.
\end{enumerate}

A modelagem adotada paralelizou somente o primeiro nível.

O teste de divisores permaneceu sequencial dentro de cada thread.

Portanto, a estrutura adotada foi:

\begin{verbatim}
parallel for
|
+-- candidato 1
|   +-- teste sequencial de divisores
|
+-- candidato 2
|   +-- teste sequencial de divisores
|
+-- ...
\end{verbatim}

Não foram utilizados paralelismo aninhado ou a cláusula
\texttt{collapse}.

Essa decisão evita a criação de regiões paralelas adicionais em uma
operação relativamente pequena e irregular e mantém o experimento alinhado
ao objetivo do enunciado: paralelizar o laço principal.


\subsection{Condição de corrida}

A contagem utiliza uma variável global à computação lógica,

\[
\texttt{count},
\]

incrementada sempre que um número primo é encontrado.

Na execução sequencial,

\begin{verbatim}
count++;
\end{verbatim}

não apresenta problema.

Entretanto, em uma execução paralela, múltiplas threads podem tentar
atualizar o mesmo contador simultaneamente.

A operação de incremento pode ser entendida conceitualmente como:

\begin{enumerate}
    \item ler o valor atual;
    \item adicionar 1;
    \item escrever o novo valor.
\end{enumerate}

Duas threads podem, por exemplo, ler simultaneamente o mesmo valor e
posteriormente sobrescrever uma atualização realizada pela outra.

Essa situação caracteriza uma \textit{race condition}, ou condição de
corrida.

Dessa forma,

\[
\texttt{parallel for}
\]

por si só não garante a correção de um programa paralelo.


\subsection{Redução}

A operação realizada sobre o contador é uma soma de contribuições
independentes.

Cada thread pode produzir uma contagem parcial,

\[
C_0,C_1,\ldots,C_{T-1},
\]

e a contagem final é

\[
C =
\sum_{i=0}^{T-1} C_i.
\]

Esse padrão corresponde diretamente à operação de redução disponível em
OpenMP.

A cláusula

\begin{verbatim}
reduction(+:count)
\end{verbatim}

permite que as threads trabalhem com contadores privados e que esses
contadores sejam combinados ao término da região paralela.

Dessa forma, a redução fornece simultaneamente paralelismo e correção
funcional.


\subsection{Distribuição de carga}

Embora cada iteração do laço externo represente o teste de um candidato,
o custo das iterações não é uniforme.

Um número par pode ser eliminado quase imediatamente, enquanto um número
primo grande pode exigir diversos testes de divisibilidade.

Portanto,

\[
W(i) \neq W(j),
\]

onde \(W(x)\) representa o trabalho necessário para determinar a
primalidade do candidato \(x\).

Esse comportamento introduz a possibilidade de desequilíbrio de carga:
uma thread pode concluir suas iterações antes das demais e permanecer
aguardando a sincronização final.


\subsection{Speedup}

A principal métrica derivada utilizada foi o speedup,

\[
S =
\frac{T_{\mathrm{seq}}}
     {T_{\mathrm{omp}}},
\]

onde:

\begin{itemize}
    \item \(T_{\mathrm{seq}}\) é o tempo da versão sequencial;
    \item \(T_{\mathrm{omp}}\) é o tempo da versão paralela correta.
\end{itemize}

Assim,

\[
S > 1
\]

indica ganho de desempenho com a paralelização, enquanto

\[
S < 1
\]

indica que o overhead paralelo foi superior ao benefício obtido.


% ============================================================
\section{Interpretação didática da tarefa}
\label{sec:atvd05-interpretacao}

A atividade não foi interpretada como um estudo de algoritmos avançados
para geração de números primos.

O problema de primalidade funciona como uma carga computacional simples,
capaz de expor conceitos fundamentais de programação paralela.

Por esse motivo, não foram utilizados métodos como o Crivo de Eratóstenes
ou algoritmos distintos entre as versões.

O núcleo computacional foi mantido igual, modificando-se apenas a forma
como as iterações do laço externo são executadas.

A sequência didática explorada foi:

\begin{verbatim}
Algoritmo sequencial
        |
        v
parallel for
        |
        v
contador compartilhado
        |
        v
race condition
        |
        v
reduction
        |
        v
programa paralelo correto
        |
        v
comparacao de desempenho
\end{verbatim}

Essa sequência evidencia um princípio fundamental de HPC: paralelizar uma
computação não consiste apenas em criar múltiplas threads. A preservação da
semântica do programa é requisito anterior à análise de desempenho.


% ============================================================
\section{Modelagem da solução}
\label{sec:atvd05-modelagem}

\subsection{Estrutura geral}

Foram modeladas três versões do algoritmo.

\subsubsection{Versão sequencial}

A primeira implementação constitui o \textit{baseline} do experimento.

Seu comportamento é:

\begin{verbatim}
para candidato = 2 ate N
    se candidato for primo
        count++
\end{verbatim}

Essa versão fornece tanto a referência funcional quanto a referência de
tempo.


\subsubsection{Versão OpenMP ingênua}

A segunda versão paraleliza diretamente o laço externo, mantendo o contador
compartilhado sem sincronização.

Conceitualmente:

\begin{verbatim}
#pragma omp parallel for

para candidato = 2 ate N
    se candidato for primo
        count++
\end{verbatim}

Essa implementação foi mantida propositalmente incorreta para demonstrar o
efeito de uma condição de corrida.

Ela participa da validação funcional, mas não do benchmark oficial de
speedup.


\subsubsection{Versão OpenMP com redução}

A terceira implementação acrescenta a redução sobre o contador:

\begin{verbatim}
#pragma omp parallel for reduction(+:count)
\end{verbatim}

Essa é a versão paralela utilizada para a comparação válida de desempenho.


\subsection{Principais funções e componentes}

A implementação foi organizada em funções com responsabilidades distintas.

\begin{center}
\begin{tabular}{|p{0.36\textwidth}|p{0.54\textwidth}|}
\hline
\textbf{Componente} & \textbf{Responsabilidade} \\
\hline
\texttt{is\_prime()} &
Verificar se um inteiro é primo. \\
\hline
\texttt{count\_primes\_sequential()} &
Realizar a contagem sequencial entre 2 e \(n\). \\
\hline
\texttt{count\_primes\_parallel\_naive()} &
Executar o laço com OpenMP mantendo a condição de corrida proposital. \\
\hline
\texttt{count\_primes\_parallel\_reduction()} &
Executar a versão paralela funcionalmente correta utilizando
\texttt{reduction}. \\
\hline
\texttt{validate\_is\_prime()} &
Validar casos conhecidos da função de primalidade. \\
\hline
\texttt{get\_actual\_thread\_count()} &
Confirmar a quantidade efetivamente criada de threads. \\
\hline
\texttt{run\_timed\_batch()} &
Executar lotes cronometrados de chamadas da função de contagem. \\
\hline
\texttt{calibrate\_internal\_repetitions()} &
Determinar automaticamente o número de repetições internas necessário
para produzir lotes mensuráveis. \\
\hline
\texttt{run\_warmups()} &
Realizar aquecimentos anteriores às medições oficiais. \\
\hline
\texttt{measure\_per\_execution()} &
Converter o tempo total do lote no tempo equivalente a uma execução. \\
\hline
\texttt{calculate\_median()} &
Calcular a mediana das vinte medições. \\
\hline
\texttt{main()} &
Coordenar configuração, validação, benchmark e apresentação dos resultados. \\
\hline
\end{tabular}
\end{center}


% ============================================================
\section{Decisões de projeto e trade-offs}
\label{sec:atvd05-decisoes}

\subsection{Preservação do algoritmo}

As versões sequencial e paralela utilizam exatamente a mesma função
\texttt{is\_prime()}.

Isso impede que eventuais diferenças de desempenho sejam atribuídas à
utilização de algoritmos distintos.

A variável experimental permanece sendo a estratégia de execução.


\subsection{Paralelização somente do laço externo}

A opção por não paralelizar o teste interno de divisores reduz o overhead e
evita a criação de paralelismo aninhado.

O laço externo apresenta unidades de trabalho independentes e constitui o
nível natural de distribuição entre as threads.


\subsection{Uso de \texttt{reduction} em vez de
\texttt{critical} ou \texttt{atomic}}

Seria possível proteger o incremento do contador por meio de mecanismos
como \texttt{critical} ou \texttt{atomic}.

Entretanto, ambos introduziriam sincronização sobre atualizações
individuais.

Como o problema consiste em somar contribuições independentes,
\texttt{reduction} representa diretamente a operação matemática necessária
e permite acumulação privada por thread.

Por esse motivo, foi adotada como solução principal.


\subsection{Escalonamento estático}

A versão paralela foi executada com

\begin{verbatim}
schedule(static)
\end{verbatim}

durante todo o experimento.

Essa decisão mantém a política de distribuição fixa e reproduzível e evita
introduzir uma segunda variável experimental.

O trade-off é que a distribuição estática de quantidades semelhantes de
iterações não garante distribuição uniforme de trabalho, pois os testes de
primalidade possuem custos distintos.


\subsection{Quatro threads fixas}

O número de threads foi fixado em quatro.

A atividade não foi transformada em um estudo de escalabilidade em função
do número de threads.

Essa decisão permite atribuir as diferenças observadas principalmente à
mudança entre execução sequencial e paralela.


\subsection{Repetições internas}

Uma primeira execução experimental mostrou que, para \(n=10\,000\), os
tempos individuais eram inferiores à resolução prática observada da
temporização, resultando em valores apresentados como zero e,
consequentemente, em um speedup indefinido.

A solução adotada foi agrupar várias execuções da mesma função em um único
lote cronometrado.

Para \(R\) repetições,

\[
T_{\mathrm{exec}}
=
\frac{T_{\mathrm{lote}}}{R}.
\]

O número de repetições é calibrado automaticamente até que as versões
sequencial e paralela produzam lotes de aproximadamente pelo menos
\(100\,\mathrm{ms}\).

Essa estratégia preserva o algoritmo original e melhora a qualidade da
medição de cargas pequenas.


\subsection{Mediana como estatística principal}

Cada configuração foi medida vinte vezes.

Foi utilizada a mediana em vez de uma única execução ou do menor tempo
observado.

Para vinte valores ordenados,

\[
t_{(1)} \leq t_{(2)} \leq \cdots \leq t_{(20)},
\]

a mediana utilizada é

\[
\widetilde{T}
=
\frac{t_{(10)}+t_{(11)}}{2}.
\]

A mediana reduz a influência de interrupções ocasionais e outros
\textit{outliers} produzidos pelo sistema operacional.


% ============================================================
\section{Representação estrutural da implementação}
\label{sec:atvd05-estrutura}

A organização lógica do programa pode ser representada da seguinte forma:

\begin{verbatim}
                         +------------------+
                         |       main       |
                         +---------+--------+
                                   |
              +--------------------+--------------------+
              |                    |                    |
              v                    v                    v
       validacao inicial     configuracao OMP      valores de N
              |                    |
              |              threads = 4
              |              dynamic = off
              |
              v
        +-----------+
        | is_prime  |
        +-----+-----+
              |
       +------+------+----------------+
       |             |                |
       v             v                v
  sequencial     OMP ingenua     OMP reduction
       |             |                |
       |        race condition        |
       |                              |
       +---------------+--------------+
                       |
                       v
                 validacao
                       |
                       v
             benchmark das versoes
                 funcionalmente
                    corretas
                       |
                       v
                  mediana
                       |
                       v
                   speedup
\end{verbatim}

A função de primalidade é compartilhada pelas três versões, evitando
diferenças algorítmicas entre as implementações.


% ============================================================
\section{Fluxo de execução}
\label{sec:atvd05-fluxo}

O fluxo experimental implementado foi:

\begin{verbatim}
INICIO
  |
  v
Configurar OpenMP
  |
  +-- desabilitar ajuste dinamico
  +-- definir 4 threads
  |
  v
Validar is_prime()
  |
  v
Confirmar numero real de threads
  |
  v
Para cada N
  |
  +-- executar sequencial
  |
  +-- executar OpenMP ingenua
  |
  +-- executar OpenMP reduction
  |
  +-- validar sequencial contra valor conhecido
  |
  +-- validar reduction contra sequencial
  |
  v
Imprimir tabela de validacao
  |
  v
Para cada N
  |
  +-- calibrar repeticoes internas
  |
  +-- executar 3 warm-ups
  |
  +-- realizar 20 medicoes
  |      |
  |      +-- alternar SEQ -> OMP
  |      +-- alternar OMP -> SEQ
  |
  +-- calcular mediana
  |
  +-- calcular speedup
  |
  +-- imprimir resultado
  |
  v
FIM
\end{verbatim}


% ============================================================
\section{Metodologia experimental}
\label{sec:atvd05-metodologia}

\subsection{Ambiente de execução}

A atividade foi executada em ambiente Windows com MSYS2 utilizando o
subsistema UCRT64.

O compilador validado foi:

\begin{itemize}
    \item GCC 15.1.0;
    \item executável \texttt{/ucrt64/bin/gcc};
    \item alvo \texttt{x86\_64-w64-mingw32};
    \item runtime OpenMP \texttt{libgomp};
    \item macro \texttt{\_OPENMP = 201511}, correspondente ao suporte
    a OpenMP 4.5.
\end{itemize}

A compilação foi realizada com:

\begin{verbatim}
gcc -std=c11 -O2 -Wall -Wextra -Wpedantic \
    -fopenmp tarefa5.c -o tarefa5.exe
\end{verbatim}

O mesmo executável contém as versões sequencial e paralela, garantindo o
mesmo compilador e o mesmo nível de otimização para todas as comparações.


\subsection{Valores de entrada}

Foram utilizados cinco valores para o limite superior:

\[
n \in
\left\{
10^4,
10^5,
10^6,
5\times10^6,
10^7
\right\}.
\]

As respectivas contagens conhecidas de números primos foram:

\begin{center}
\begin{tabular}{|r|r|}
\hline
\textbf{\(n\)} & \textbf{Quantidade de primos} \\
\hline
10\,000 & 1\,229 \\
100\,000 & 9\,592 \\
1\,000\,000 & 78\,498 \\
5\,000\,000 & 348\,513 \\
10\,000\,000 & 664\,579 \\
\hline
\end{tabular}
\end{center}


\subsection{Configuração das medições}

A configuração experimental utilizada foi:

\begin{center}
\begin{tabular}{|l|l|}
\hline
\textbf{Parâmetro} & \textbf{Configuração} \\
\hline
Threads OpenMP & 4 \\
\hline
Escalonamento & \texttt{static} \\
\hline
Aquecimentos & 3 \\
\hline
Medições & 20 \\
\hline
Estatística & mediana \\
\hline
Relógio & \texttt{omp\_get\_wtime()} \\
\hline
Lote mínimo de temporização & aproximadamente 100 ms \\
\hline
Otimização & \texttt{-O2} \\
\hline
\end{tabular}
\end{center}


\subsection{Região cronometrada}

Somente a operação de contagem de primos foi incluída na região
cronometrada.

A estrutura utilizada foi equivalente a:

\begin{verbatim}
inicio = omp_get_wtime();

executar contagem;

fim = omp_get_wtime();
\end{verbatim}

Na versão OpenMP, isso inclui o custo efetivo de:

\begin{itemize}
    \item entrada na região paralela;
    \item distribuição das iterações;
    \item execução das threads;
    \item redução dos resultados;
    \item sincronização final.
\end{itemize}

Foram mantidos fora da região medida:

\begin{itemize}
    \item impressão dos resultados;
    \item cálculo da mediana;
    \item cálculo do speedup;
    \item comparação funcional;
    \item configuração do OpenMP.
\end{itemize}


\subsection{Aquecimentos}

Antes das medições oficiais foram executados três aquecimentos por
configuração.

Essa etapa reduz a influência da primeira ativação do runtime OpenMP e de
efeitos transitórios da execução inicial.


\subsection{Alternância da ordem}

Para reduzir viés sistemático de ordem, as versões foram alternadas entre
as medições.

O padrão utilizado foi:

\begin{verbatim}
medicao 1:  SEQ -> OMP
medicao 2:  OMP -> SEQ
medicao 3:  SEQ -> OMP
medicao 4:  OMP -> SEQ
...
\end{verbatim}

Dessa forma, nenhuma implementação é sistematicamente beneficiada ou
prejudicada por ser executada sempre primeiro.


\subsection{Calibração das repetições internas}

O número de repetições internas foi determinado automaticamente por
duplicação:

\[
1,2,4,8,16,\ldots
\]

até que os lotes das implementações válidas alcançassem a duração mínima
estabelecida.

A mesma quantidade de repetições foi utilizada para as duas versões dentro
de cada valor de \(n\).

Os tempos apresentados correspondem ao tempo normalizado de uma única
contagem.


\subsection{Prevenção de eliminação da computação}

O resultado da computação possui efeito observável e é posteriormente
validado.

Além disso, o \textit{harness} de benchmark utiliza um valor protegido para
o limite \(n\) durante as repetições internas.

Não foi utilizado \texttt{volatile} no algoritmo de primalidade propriamente
dito, evitando alterar artificialmente seu comportamento de otimização.


% ============================================================
\section{Código-fonte desenvolvido}
\label{sec:atvd05-codigo}

\lstinputlisting[
    style=codigoC,
    caption={Código-fonte da Atividade 5.},
    label={lst:atvd05-codigo}
]{../ATVD5/tarefa5.c}


% ============================================================
\section{Validação da implementação}
\label{sec:atvd05-validacao}

A validação foi realizada em três níveis.


\subsection{Validação da função de primalidade}

A função \texttt{is\_prime()} foi testada utilizando casos conhecidos,
incluindo números primos, compostos e valores inferiores a 2.

Entre os casos utilizados estavam:

\[
0,\ 1,\ 2,\ 3,\ 4,\ 5,\ 9,\ 17,\ 25.
\]


\subsection{Validação por contagens conhecidas}

A contagem sequencial foi comparada com valores conhecidos de
\(\pi(n)\), onde \(\pi(n)\) representa a quantidade de primos menores ou
iguais a \(n\).

Todas as cinco entradas utilizadas no experimento produziram os valores
esperados.


\subsection{Comparação entre as versões}

Os resultados obtidos foram:

\begin{center}
\small
\begin{tabular}{|r|r|r|r|l|}
\hline
\textbf{\(N\)} &
\textbf{Sequencial} &
\textbf{OMP ingênua} &
\textbf{OMP reduction} &
\textbf{Status ingênua} \\
\hline
10\,000
& 1\,229
& 279
& 1\,229
& incorreta \\
\hline

100\,000
& 9\,592
& 2\,199
& 9\,592
& incorreta \\
\hline

1\,000\,000
& 78\,498
& 18\,260
& 78\,498
& incorreta \\
\hline

5\,000\,000
& 348\,513
& 81\,795
& 348\,513
& incorreta \\
\hline

10\,000\,000
& 664\,579
& 156\,318
& 664\,579
& incorreta \\
\hline
\end{tabular}
\end{center}

A versão sequencial e a versão com redução coincidiram em todos os casos:

\[
C_{\mathrm{seq}}(n)
=
C_{\mathrm{reduction}}(n).
\]

A implementação ingênua apresentou contagens inferiores em todos os
tamanhos testados.

Isso confirma experimentalmente a presença da condição de corrida.

Os valores produzidos pela versão ingênua não devem ser interpretados como
uma fração determinística da resposta correta. Uma condição de corrida não
oferece garantia sobre o valor final, mesmo que resultados semelhantes
possam ocorrer em execuções realizadas sob as mesmas condições.


% ============================================================
\section{Resultados experimentais}
\label{sec:atvd05-resultados}

Após a calibração temporal e a execução das vinte medições por configuração,
foram obtidos os seguintes resultados finais.

\begin{center}
\small
\begin{tabular}{|r|r|r|r|r|r|}
\hline
\textbf{\(N\)} &
\textbf{Primos} &
\textbf{Reps} &
\textbf{Seq. (ms)} &
\textbf{OpenMP (ms)} &
\textbf{Speedup} \\
\hline

10\,000
& 1\,229
& 512
& 0,282
& 0,195
& 1,45$\times$ \\
\hline

100\,000
& 9\,592
& 64
& 5,836
& 2,984
& 1,96$\times$ \\
\hline

1\,000\,000
& 78\,498
& 2
& 170,250
& 72,500
& 2,35$\times$ \\
\hline

5\,000\,000
& 348\,513
& 1
& 1\,262,000
& 603,000
& 2,09$\times$ \\
\hline

10\,000\,000
& 664\,579
& 1
& 3\,397,000
& 1\,610,000
& 2,11$\times$ \\
\hline
\end{tabular}
\end{center}

Os valores de \texttt{Reps} correspondem ao número de execuções internas
utilizadas em cada lote de temporização.

O valor reportado nas colunas de tempo continua correspondendo a uma única
execução, pois o tempo total do lote foi normalizado pelo número de
repetições.


% ============================================================
\section{Visualização dos resultados}
\label{sec:atvd05-visualizacao}

A evolução do speedup pode ser resumida por:

\begin{center}
\begin{tabular}{|r|c|}
\hline
\textbf{\(N\)} & \textbf{Speedup} \\
\hline
10\,000 & 1,45$\times$ \\
100\,000 & 1,96$\times$ \\
1\,000\,000 & 2,35$\times$ \\
5\,000\,000 & 2,09$\times$ \\
10\,000\,000 & 2,11$\times$ \\
\hline
\end{tabular}
\end{center}

Uma representação textual da tendência observada é:

\begin{verbatim}
Speedup

2.5 |                     *
    |
2.0 |          *                    *          *
    |
1.5 |    *
    |
1.0 +------------------------------------------------
        10^4     10^5      10^6      5x10^6     10^7

                         N
\end{verbatim}

O maior speedup da execução final foi observado em

\[
N=1\,000\,000,
\]

com

\[
S=2,35\times.
\]

Nos dois maiores tamanhos avaliados, os valores observados ficaram próximos
de

\[
2,1\times.
\]


% ============================================================
\section{Análise e discussão}
\label{sec:atvd05-discussao}

\subsection{Correção da versão sequencial}

A implementação sequencial produziu exatamente as contagens conhecidas para
todos os valores de entrada avaliados.

Isso estabelece uma referência funcional confiável para a comparação com
as implementações paralelas.


\subsection{Falha da paralelização ingênua}

A inclusão de

\begin{verbatim}
#pragma omp parallel for
\end{verbatim}

sem tratamento adicional do contador compartilhado produziu resultados
incorretos em todos os tamanhos avaliados.

Por exemplo, para

\[
N=10\,000\,000,
\]

a resposta correta foi

\[
664\,579,
\]

enquanto a versão ingênua retornou

\[
156\,318.
\]

A diferença não decorre da função de primalidade, pois o mesmo
\texttt{is\_prime()} é utilizado nas três versões.

A causa está na atualização concorrente de \texttt{count}.

Dessa forma, o experimento demonstra concretamente que:

\[
\boxed{
\text{paralelização sintaticamente válida}
\not\Rightarrow
\text{programa funcionalmente correto}
}
\]


\subsection{Correção por redução}

A versão com

\begin{verbatim}
reduction(+:count)
\end{verbatim}

reproduziu exatamente a contagem sequencial para todas as entradas.

Isso ocorre porque OpenMP mantém contribuições parciais das threads e as
combina por soma ao final da região paralela.

Portanto, a redução corresponde diretamente à estrutura matemática do
problema.


\subsection{Efeito do tamanho do problema}

Para \(N=10\,000\), o speedup foi:

\[
S=1,45\times.
\]

Embora a versão paralela já tenha sido mais rápida, o benefício foi menor
que nos tamanhos intermediários.

Isso é coerente com o efeito do overhead paralelo.

O tempo da implementação OpenMP pode ser representado conceitualmente por:

\[
T_{\mathrm{omp}}
=
T_{\mathrm{computacao}}
+
T_{\mathrm{overhead}},
\]

onde o segundo termo inclui custos de gerenciamento da região paralela,
distribuição das iterações, redução e sincronização.

Quando o problema é pequeno, esse custo representa parcela proporcionalmente
maior do tempo total.


\subsection{Speedup observado}

Os speedups obtidos foram:

\[
1,45\times,\quad
1,96\times,\quad
2,35\times,\quad
2,09\times,\quad
2,11\times.
\]

A melhor aceleração observada ocorreu para

\[
N=10^6,
\]

com

\[
S=2,35\times.
\]

Nos dois maiores tamanhos, os ganhos ficaram próximos de
\(2,1\times\).

A utilização de quatro threads não implica um speedup de \(4\times\).

Para atingir o speedup ideal seria necessário, entre outros fatores:

\begin{itemize}
    \item ausência de overhead de paralelização;
    \item trabalho perfeitamente balanceado;
    \item ausência de sincronização;
    \item disponibilidade integral dos recursos computacionais.
\end{itemize}

Essas condições não ocorrem no experimento real.


\subsection{Distribuição de carga}

A política utilizada foi \texttt{schedule(static)}.

Ela distribui previamente conjuntos de iterações entre as threads e possui
baixo overhead.

Entretanto, igual quantidade de candidatos não corresponde necessariamente
a igual quantidade de trabalho.

Um número par pode ser rejeitado rapidamente, enquanto um primo grande pode
exigir testes de diversos divisores.

Além disso, candidatos maiores apresentam um limite potencial de teste
também maior.

Assim,

\[
\text{mesmo número de iterações}
\not\Rightarrow
\text{mesma carga computacional}.
\]

Consequentemente, algumas threads podem terminar seus blocos antes de outras
e permanecer aguardando a barreira implícita do final do
\texttt{parallel for}.

Esse desequilíbrio contribui para limitar o speedup observado.


\subsection{Crescimento do custo sequencial}

Os tempos sequenciais cresceram de

\[
0,282\,\mathrm{ms}
\]

para

\[
3397\,\mathrm{ms}
\]

ao longo da faixa avaliada.

O crescimento não é linear em \(N\), pois o custo de testar um candidato
também cresce com seu valor e depende de seus fatores.

Dessa forma, aumentar \(N\) aumenta simultaneamente:

\begin{enumerate}
    \item a quantidade de candidatos avaliados;
    \item o custo potencial de algumas avaliações individuais.
\end{enumerate}


\subsection{Efeito da calibração temporal}

A calibração das repetições internas mostrou-se necessária para os menores
valores de \(N\).

As quantidades finais foram:

\[
512,\ 64,\ 2,\ 1,\ 1.
\]

Isso produz o comportamento esperado:

\[
N \uparrow
\quad\Rightarrow\quad
T_{\mathrm{execucao}} \uparrow
\quad\Rightarrow\quad
R \downarrow.
\]

Para os dois maiores tamanhos, uma única contagem já possui duração
suficiente para a medição.

A técnica eliminou o problema inicialmente observado de tempos reportados
como zero para cargas pequenas sem alterar o algoritmo que estava sendo
comparado.


% ============================================================
\section{Limitações e ameaças à validade}
\label{sec:atvd05-limitacoes}

Os resultados devem ser interpretados considerando algumas limitações.

\subsection{Uma única máquina}

As medições correspondem ao ambiente utilizado na atividade.

Arquiteturas diferentes podem apresentar outros níveis de speedup em razão
de diferenças no número de núcleos, memória, frequência, sistema operacional
e implementação do runtime OpenMP.


\subsection{Número fixo de threads}

Foram utilizadas quatro threads em todas as execuções.

Portanto, a atividade não permite concluir sobre escalabilidade para
quantidades maiores ou menores de threads.


\subsection{Política de escalonamento fixa}

Apenas \texttt{schedule(static)} foi utilizado.

Outras políticas de OpenMP podem responder diferentemente ao custo
irregular das iterações.

A comparação entre políticas de escalonamento não fazia parte do escopo
experimental desta atividade.


\subsection{Interferência do sistema operacional}

Mesmo com aquecimentos e múltiplas medições, processos externos,
escalonamento do sistema operacional, alterações de frequência da CPU e
condições térmicas podem introduzir variação temporal.

A utilização da mediana reduz, mas não elimina completamente, esse efeito.


\subsection{Ausência de afinidade explícita}

Não foi implementado controle explícito de afinidade de threads aos
processadores.

A localização efetiva das threads foi, portanto, determinada pelo runtime e
pelo sistema operacional.


\subsection{Versão ingênua}

A versão com condição de corrida foi mantida exclusivamente para
demonstração.

Seus resultados não podem ser utilizados para avaliação de desempenho
válida, pois ela não preserva a semântica do programa sequencial.

Uma implementação incorreta não pode ser considerada uma otimização mesmo
quando executa em menor tempo.


\subsection{Resolução temporal}

A primeira modelagem baseada em uma única contagem por medição mostrou-se
insuficiente para os menores valores de \(N\).

A utilização posterior de lotes calibrados mitigou essa limitação, mas
também tornou necessário normalizar o tempo pela quantidade de repetições.


% ============================================================
\section{Conclusões da atividade}
\label{sec:atvd05-conclusao}

A atividade demonstrou que o laço externo da contagem de números primos é
adequado ao paralelismo de dados, pois os diferentes candidatos podem ser
avaliados independentemente.

Entretanto, a simples aplicação de

\begin{verbatim}
#pragma omp parallel for
\end{verbatim}

não foi suficiente para preservar a correção.

A atualização concorrente do contador introduziu uma condição de corrida,
produzindo resultados incorretos para todos os tamanhos avaliados.

A utilização de

\begin{verbatim}
reduction(+:count)
\end{verbatim}

eliminou esse problema e fez com que a implementação OpenMP reproduzisse
exatamente os resultados da versão sequencial.

Após a validação funcional, a comparação de desempenho mostrou benefício
consistente da paralelização.

Os speedups observados variaram entre

\[
1,45\times
\]

e

\[
2,35\times.
\]

O melhor resultado foi obtido para \(N=10^6\), enquanto os dois maiores
problemas apresentaram ganhos próximos de \(2,1\times\).

A atividade também evidenciou que o desempenho paralelo depende não apenas
da quantidade de threads, mas de fatores como overhead, granularidade,
sincronização e balanceamento de carga.

Finalmente, a necessidade de calibração temporal para cargas pequenas
demonstrou que a metodologia de benchmark também constitui parte essencial
de uma avaliação de desempenho: medições rápidas demais podem ser dominadas
pela resolução do relógio e produzir conclusões inadequadas.


% ============================================================
\section{Síntese para defesa oral}
\label{sec:atvd05-defesa}

Os principais pontos para apresentação oral da atividade são:

\begin{enumerate}
    \item \textbf{Problema:}
    contar quantos números primos existem entre 2 e \(N\).

    \item \textbf{Algoritmo-base:}
    o mesmo teste de primalidade foi utilizado em todas as versões para
    garantir uma comparação justa.

    \item \textbf{Paralelização:}
    somente o laço externo dos candidatos foi paralelizado.
    O laço interno dos divisores permaneceu sequencial.

    \item \textbf{Versões:}
    foram implementadas uma versão sequencial, uma versão OpenMP ingênua e
    uma versão OpenMP com \texttt{reduction}.

    \item \textbf{Problema da versão ingênua:}
    múltiplas threads atualizam \texttt{count++} simultaneamente,
    produzindo uma condição de corrida.

    \item \textbf{Evidência da condição de corrida:}
    para \(N=10^7\), a contagem correta foi \(664\,579\), enquanto a versão
    ingênua retornou \(156\,318\).

    \item \textbf{Correção:}
    a versão com \texttt{reduction} coincidiu com a versão sequencial para
    todos os cinco valores de \(N\).

    \item \textbf{Threads:}
    foram utilizadas quatro threads fixas, com ajuste dinâmico desabilitado.

    \item \textbf{Scheduling:}
    foi utilizado \texttt{schedule(static)} em toda a atividade para manter
    essa variável controlada.

    \item \textbf{Temporização:}
    foi utilizado \texttt{omp\_get\_wtime()}.

    \item \textbf{Metodologia:}
    foram utilizados 3 warm-ups, 20 medições e a mediana como estatística
    principal.

    \item \textbf{Problema de resolução temporal:}
    execuções pequenas inicialmente produziram tempos iguais a zero.
    O problema foi corrigido por repetição interna calibrada até lotes de
    aproximadamente 100 ms.

    \item \textbf{Speedup:}
    foi calculado por
    \[
    S=\frac{T_{\mathrm{seq}}}{T_{\mathrm{omp}}}.
    \]

    \item \textbf{Resultado de desempenho:}
    os speedups finais foram
    \[
    1,45\times,\;
    1,96\times,\;
    2,35\times,\;
    2,09\times,\;
    2,11\times.
    \]

    \item \textbf{Melhor speedup observado:}
    \[
    2,35\times
    \]
    para \(N=1\,000\,000\).

    \item \textbf{Por que quatro threads não geraram \(4\times\):}
    existem overhead de OpenMP, sincronização, competição por recursos e
    desequilíbrio de carga.

    \item \textbf{Distribuição de carga:}
    candidatos diferentes possuem custos diferentes; números compostos
    simples podem terminar rapidamente, enquanto primos grandes demandam
    mais testes.

    \item \textbf{Conclusão principal:}
    primeiro é necessário garantir correção; somente depois faz sentido
    avaliar desempenho. A versão com \texttt{reduction} foi simultaneamente
    correta e mais rápida que a implementação sequencial nos tamanhos
    avaliados.
\end{enumerate}